\documentclass[11pt,a4paper]{article}

% ========================
% ESSENTIAL PACKAGES
% ========================

% Language and encoding
\usepackage[utf8]{inputenc}
\usepackage[T1]{fontenc}
\usepackage[english]{babel}
%\usepackage{natbib}

% Math packages
\usepackage{amsmath}        % Enhanced math environments
\usepackage{amssymb}        % Additional math symbols
\usepackage{amsthm}         % Theorem environments
\usepackage{amsfonts}       % Additional math fonts
\usepackage{mathtools}      % Extension to amsmath
\usepackage{bm}             % Bold math symbols
\usepackage{dsfont}         % Double-struck fonts (for sets like ℝ, ℕ)
\usepackage{mathrsfs}       % Script fonts for math

% Graphics and figures
\usepackage{graphicx}       % Include graphics
\usepackage{float}          % Better float placement
\usepackage{subcaption}     % Subfigures
\usepackage{tikz}           % Drawing diagrams
\usepackage{pgfplots}       % Plotting
\pgfplotsset{compat=1.18}

% Page layout and formatting
\usepackage[margin=1in]{geometry}
\usepackage{fancyhdr}       % Custom headers and footers
\usepackage{titlesec}       % Customize section titles
\usepackage{enumitem}       % Customize lists
\usepackage{parskip}        % Paragraph spacing

% Colors and highlighting
\usepackage{xcolor}
\usepackage{mdframed}       % Framed environments
\usepackage{tcolorbox}      % Colored boxes

% Tables
\usepackage{array}          % Extended array environments
\usepackage{booktabs}       % Professional tables
\usepackage{multirow}       % Multi-row cells

% References and links
\usepackage{hyperref}       % Hyperlinks
\usepackage{cleveref}       % Smart cross-references

% Additional useful packages
\usepackage{cancel}         % Cancel terms in equations
\usepackage{siunitx}        % SI units
\usepackage{listings}       % Code listings (if needed)

% ========================
% CUSTOM COMMANDS
% ========================

% Common math sets
\newcommand{\R}{\mathbb{R}}
\newcommand{\N}{\mathbb{N}}
\newcommand{\Z}{\mathbb{Z}}
\newcommand{\Q}{\mathbb{Q}}
\newcommand{\C}{\mathbb{C}}

% Common operators
\DeclareMathOperator{\lcm}{lcm}
\DeclareMathOperator{\Tr}{Tr}
\DeclareMathOperator{\rank}{rank}
%\DeclareMathOperator{\span}{span}
%\DeclareMathOperator{\null}{null}

% Shortcuts for common expressions
\newcommand{\abs}[1]{\left|#1\right|}
\newcommand{\norm}[1]{\left\|#1\right\|}
\newcommand{\inner}[2]{\left\langle #1, #2 \right\rangle}
\newcommand{\ceil}[1]{\left\lceil #1 \right\rceil}
\newcommand{\floor}[1]{\left\lfloor #1 \right\rfloor}

% Derivatives and integrals
\newcommand{\dd}[1]{\,\mathrm{d}#1}
\newcommand{\dv}[2]{\frac{\mathrm{d}#1}{\mathrm{d}#2}}
\newcommand{\pdv}[2]{\frac{\partial #1}{\partial #2}}

\renewcommand{\hat}{\widehat}
\renewcommand{\tilde}{\widetilde}

% ========================
% THEOREM ENVIRONMENTS
% ========================

\theoremstyle{definition}
\newtheorem{definition}{Definition}[section]
\newtheorem{example}{Example}[section]
\newtheorem{exercise}{Exercise}[section]

\theoremstyle{plain}
\newtheorem{theorem}{Theorem}[section]
\newtheorem{lemma}{Lemma}[section]
\newtheorem{corollary}{Corollary}[section]
\newtheorem{proposition}{Proposition}[section]
\newtheorem{assumption}{Assumptio}[section]

\theoremstyle{remark}
\newtheorem{remark}{Remark}[section]
\newtheorem{note}{Note}[section]

\newcommand{\supp}{\mathrm{supp}}
\newcommand{\op}{\mathrm{op}}

\newcommand{\mX}{\mathcal X}
\newcommand{\mD}{\mathcal D}
\newcommand{\bP}{\mathbb P}
\newcommand{\lp}{L^2(\bP)}
\newcommand{\lx}{L^2(\mX)}
\newcommand{\ltwo}{L^2(\mX)}
\newcommand{\ud}{\mathrm d}
\newcommand{\mH}{\mathcal H}
\newcommand{\mA}{\mathcal A}

\newcommand{\diag}{\mathrm{diag}}

% ========================
% COLORED BOXES
% ========================

% Important theorem box
\newtcolorbox{importantbox}{
	colback=blue!5!white,
	colframe=blue!75!black,
	title=Important
}

% Warning/Note box
\newtcolorbox{notebox}{
	colback=yellow!10!white,
	colframe=orange!75!black,
	title=Note
}

% Example box
\newtcolorbox{examplebox}{
	colback=green!5!white,
	colframe=green!75!black,
	title=Example
}

% ========================
% DOCUMENT SETUP
% ========================

% Page style
\pagestyle{fancy}
\fancyhf{}
\rhead{\thepage}
\lhead{\leftmark}

% Title formatting
\title{Functional Bandit Notes}
\author{Your Name}
\date{\today}

% Hyperlink setup
\hypersetup{
	colorlinks=true,
	linkcolor=blue,
	filecolor=magenta,      
	urlcolor=cyan,
	citecolor=red
}

\begin{document}

The main purpose of this document is to provide some insights into minimax rates in functional regression, from a lower bound perspective.

For simplicity, assume we use the identity kernel. Let $\{\psi_i\}_{i=1}^{\infty}$ be an orthonormal basis of $L_2([0,1])$ that are eigenfunctions of the population covariance operator.
Define the following two function classes:
\begin{align*}
\textstyle
\Theta_2 &:= \{\beta\in L_2([0,1]): \sum_{j=1}^{\infty}\vartheta_j^2 =\sum_{j=1}^{\infty}\langle \beta,\psi_j\rangle^2\leq 1\};\\
\Theta_1 &:= \{\beta\in\Theta_2: \sum_{j=1}^{\infty}|\vartheta_j| = \sum_{j=1}^{\infty}|\langle\beta,\psi_j\rangle| \leq 1\}.
\end{align*}

Let $P$ be a known distribution supported on $L_2([0,1])$ such that
$$
s_j := \mathbb E_P[\langle f,\psi_j\rangle^2] \asymp j^{-2r},\;\;\;\;\;\;\forall j\geq 1.
$$
In all our lower bound constructions, the distribution $P$ satisfies the following additional properties, though some or all of them may not be necessary for certain upper bound:
\begin{enumerate}
\item For $f\sim P$, $|\langle f,\psi_j\rangle| \lesssim j^{-2r}$ for all $j\geq 1$ almost surely;
\item $s_j-s_{j+1}\gtrsim j^{-2r-1}$.
\end{enumerate}

Data are collected as $\mathcal D\{(f_i,y_i)\}_{i=1}^n$ where $f_1,\cdots,f_n\overset{i.i.d.}{\sim} P$ and $\mathbb E[y_i]=\langle f_i,\beta_0\rangle$.
For $\Theta\in\{\Theta_1,\Theta_2\}$, 
the first minimax rate we consider is the standard mean-square error:
$$
\mathfrak T_{mse}(\Theta) := \inf_{\hat\beta}\sup_{\beta_0\in\Theta} \mathbb E_{\beta_0}\left[\int \langle f,\hat\beta-\beta_0\rangle^2\ud P(f)\right],
$$
where $\mathbb E_{\beta_0}$ is over randomness in $\mathcal D$ under $\beta_0$.
The second minimax rate we consider is the sup-norm error:
$$
\mathfrak T_{sup}(\Theta) := \inf_{\hat\beta}\sup_{\beta_0\in\Theta}\mathbb E_{\beta_0}\left[\sup_{f\in\mathrm{supp}(P)}\langle f,\hat\beta-\beta_0\rangle^2\right].
$$
Next, we say a map $\varphi:f\mapsto [a,b]$, itself being a random variable depending on $\mathcal D$, is an \emph{honest} confidence interval if it satisfies the following:
$$
\Pr_{\beta_0}\left[f\in \varphi(f)\right] \geq 0.95,\;\;\;\;\;\;\forall\beta_0\in\Theta,\;\;f\in\mathrm{supp}(P),
$$
where the randomness is over $\mathcal D$ under $\beta_0$, and $\varphi$ is a random variable depending on $\mathcal D$.
Let $\Phi$ be the set of all honest confidence intervals. 
The following minimax rate is concerning the MSE of honest confidence intervals:
$$
\mathfrak T_{CI}(\Theta) := \inf_{\varphi\in\Phi}\sup_{\beta_0\in\Theta}\mathbb E_{\beta_0}\left[\int \mu(\varphi(f))^2\ud P(f)\right],
$$
where $\mu(\varphi(f))$ is the Lesbesgue measure of $\varphi(f)$.

Table \ref{tab:summary} summarizes minimax lower bounds for different cases.
\begin{table}[h]
\centering
\caption{Summary of minimax lower bounds.}
\label{tab:summary}
\begin{tabular}{cccl}
\hline
Risk& Class& Lower bound& Matching upper bound and method\\ 
\hline
$\mathfrak T_{mse}$& $\Theta_2$& $n^{-2r/(2r+1)}$& KRR with $\lambda\asymp n^{-2r/(2r+1)}$\\
$\mathfrak T_{sup}$& $\Theta_2$& $n^{-(2r-1)/(2r+1)}$& Unknown. Possibly KRR with suprema analysis\\
$\mathfrak T_{CI}$& $\Theta_2$& $n^{-(2r-1)/2r}$& KRR with $\lambda = o(n^{-1})$\\
$\mathfrak T_{mse}$& $\Theta_1$& $n^{-(2r+1)/(2r+2)}$& Unknown. Possibly trend filtering\\
$\mathfrak T_{sup}$& $\Theta_1$& $n^{-r/(r+1)}$& Unknown. Possibly trend filtering\\
$\mathfrak T_{CI}$& $\Theta_1$& $n^{-2r/(2r+1)}$& Anisotropic KRR\\
\hline
\end{tabular}
\end{table}

\subsection{Lower bounds on $\mathfrak T_{mse}$}

Using construction and techniques very similar to \cite{cai2012minimax}.
Let $\Theta\in\Theta_2$ be such that $\theta_M=\theta_{M+1}=\cdots=\theta_{2M}=\pm 1/\sqrt{M}$. (The function is constructed as $\beta = \sum_{j=M}^{2M}\theta_j\psi_j$.)
Let $\tilde\Theta\subseteq\Theta$ be selected such that $\ln|\tilde\Theta|\asymp M$ and for every $\theta,\theta'\in\Theta$, $\|\theta-\theta'\|_2\gtrsim \sqrt{M}$.
For any distinct $\theta,\theta'\in\tilde\Theta$, we have
\begin{align*}
d_2(\theta,\theta') &= \int \langle f,\theta-\theta'\rangle^2\ud P(f)\gtrsim \sum_{j=M}^{2M} \frac{1}{M}\mathbb E[\langle f,\psi_j\rangle^2] \gtrsim \sum_{j=M}^{2M} \frac{j^{-2r}}{M} = M^{-2r}.
\end{align*}
On the other hand,
\begin{align*}
\mathrm{KL}(P_{\theta}\|P_{\theta'}) \lesssim n\mathbb E_P[\langle f,\theta-\theta'\rangle^2] \lesssim n\sum_{j=M}^{2M}\frac{1}{M}\cdot \mathbb E[\langle f,\psi_j\rangle^2]\lesssim \frac{n}{M}\sum_{j=M}^{2M}j^{-2r} = nM^{-2r}.
\end{align*}
Using Fano's equality, we can set $nM^{-2r}\ll \ln|\tilde\Theta|\asymp M$, meaning $M\asymp n^{1/(2r+1)}$. This leads to
$$
\mathfrak T_{mse}(\Theta_2) \gtrsim n^{-2r/(2r+1)}.
$$

For the case of $\Theta_1$, we must set $\theta_M=\cdots\theta_{2M}=\pm 1/M$ in order to make sure $\beta\in\Theta_1$. Using the same line of calculations, we have
\begin{align*}
d_2(\theta,\theta') &\gtrsim \sum_{j=M}^{2M} \frac{j^{-2r}}{M^2} = M^{-2r-1};\\
\mathrm{KL}(Q_{\theta}\|Q_{\theta'}) &\lesssim n\sum_{j=M}^{2M}\frac{1}{M^2} j^{-2r} = nM^{-2r-1}.
\end{align*}
Equating $nM^{2r-1}$ with $M$ again we have $M=n^{1/(2r+2)}$. Hence
$$
\mathfrak T_{mse}(\Theta_1) \gtrsim n^{-(2r+1)/(2r+2)}.
$$

\subsection{Lower bounds on $\mathfrak T_{sup}$}

We are going to replace $d_2$ in the above section with the following sup-norm metric (for $\Theta_2$ first):
\begin{align*}
d_{\infty}(\theta,\theta') &= \sup_{f\in\mathrm{supp}(P)}\langle f,\theta-\theta'\rangle^2 \gtrsim \sup_{f\in\mathrm{supp}(P)}\left(\sum_{j=M}^{2M}\frac{\langle f,\psi_j\rangle}{\sqrt{M}}\right)^2 \gtrsim \left(\sum_{j=M}^{2M}\frac{j^{-r}}{\sqrt{M}}\right)^2 \gtrsim M^{1-2r},
\end{align*}
where the second inequality is attained by considering $f\in\mathrm{supp}(P)$ such that $\mathrm{sgn}(\langle f,\psi_j\rangle) = \mathrm{sgn}(\theta_j)$.
The KL part is the same, yielding $M\asymp n^{1/(1+2r)}$. Subsequently, 
$$
\mathfrak T_{sup}(\Theta_2) \gtrsim n^{-(2r-1)/(2r+1)}.
$$

For the case of $\Theta_1$, the sup-norm metric is calculated as 
\begin{align*}
d_{\infty}(\theta,\theta') &= \sup_{f\in\mathrm{supp}(P)}\langle f,\theta-\theta'\rangle^2 \gtrsim \sup_{f\in\mathrm{supp}(P)}\left(\sum_{j=M}^{2M}\frac{\langle f,\psi_j\rangle}{{M}}\right)^2 \gtrsim \left(\sum_{j=M}^{2M}\frac{j^{-r}}{{M}}\right)^2 \gtrsim M^{-2r}.
\end{align*}
Using $M\asymp n^{1/(2r+2)}$, we obtain
$$
\mathfrak T_{sup}(\Theta_2) \gtrsim n^{-(2r)/(2r+2)} = n^{-r/(r+1)}.
$$

\subsection{Lower bounds on $\mathfrak T_{CI}$}

We are going to use arguments similar to \cite{cai2006adaptive}. Because $\varphi$ is honest, we have
$$
\Pr_0[0\in\varphi(f)] \geq 0.95, \;\;\;\;\;\;\forall f\in\mathrm{supp}(P),
$$
where the randomness is over $\mathcal D$ under the zero function $\beta_0=0$.

Now consider an arbitrary $f\in\mathrm{supp}(P)$ and its corresponding $\theta(f)\in\Theta_2$, such that $\mathrm{sgn}(\theta(f)_j)=\mathrm{sgn}(\langle f,\psi_j\rangle)$.
Let $\beta(f) =\sum_{j=M}^{2M}\theta(f)_j\psi_j$.
It is easy to verify that
$$
a(f) := \langle f, \beta(f)\rangle \asymp \sum_{j=M}^{2M}\frac{j^{-r}}{\sqrt{M}}\asymp M^{-r+1/2}.
$$
Because $\varphi$ is honest, we have
$$
\Pr_{\beta(f)}[a\in\varphi(f)] \geq 0.95.
$$

Now, we upper bound the TV distance between $Q_0$ and $Q_{\theta(f)}$ as follows:
\begin{align*}
\mathrm{KL}(Q_0\|Q_{\beta(f)}) \lesssim nM^{-2r}.
\end{align*}
Equating the right-hand side of the above inequality with a small positive constant, meaning that $M\asymp n^{1/2r}$, by Pinskier's inequality we have
$$
\Pr_{0}[a\in\varphi(f)] \geq 0.9.
$$
Therefore,
$$
\Pr_{0}\left([0,a]\in\varphi(f)\right) = \Omega(1) \;\;\Longrightarrow \mathbb E_0\left[\mu(\varphi(f))^2\right] \gtrsim a^2 \gtrsim M^{-2r+1}.
$$
Because the above inequality holds for any $f\in\mathrm{supp}(P)$, we have
$$
\mathfrak T_{CI}(\Theta_2) \gtrsim n^{-(2r-1)/2r}. 
$$

For the case of $\Theta=\Theta_1$, simply replacing the asymptotic scalings on $a$ the KL divergence with
$$
a \asymp \sum_{j=M}^{2M}\frac{j^{-r}}{{M}}\asymp M^{-r}, \;\;\;\;\;\;\mathrm{KL}(Q_0\|Q_{\beta(f)}) \lesssim n M^{-2r-1}.
$$
Setting $M\asymp n^{1/(2r+1)}$, we obtain
$$
\mathfrak T_{CI}(\Theta_2) \gtrsim n^{-2r/(2r+1)}.
$$

\bibliographystyle{plain}
\bibliography{refs}

\end{document}